\par\Needspace{8\baselineskip}\bigskip
\pdfbookmark[3]{Darstellung, 8}{pallas-II-article8}
\begin{center}8.\end{center}
\noindent\hspace*{\parindent}Bei dieser Untersuchung verwenden wir dieselben Zeichen, die wir im ersten Abschnitt eingeführt haben, sowohl für die eigentlichen Elemente als auch für die übrigen Hilfsgrößen. So bezeichnet etwa \(A\) den nun veränderlichen Wert von \(\dfrac{y\mathrm{d}z-z\mathrm{d}y}{\mathrm{d}t}\). Man sieht, dass alle endlichen Gleichungen zwischen diesen Größen \(A,B,C,F,G,H,v,E\) usw., den Elementen und den Koordinaten \(x,y,z,r\) samt ihren Differenzialen \(\dfrac{\mathrm{d}x}{\mathrm{d}t}\mathpunct{\raisebox{.45ex}{,}}\dfrac{\mathrm{d}y}{\mathrm{d}t}\mathpunct{\raisebox{.45ex}{,}}\dfrac{\mathrm{d}z}{\mathrm{d}t}\mathpunct{\raisebox{.45ex}{,}}\dfrac{\mathrm{d}r}{\mathrm{d}t}\) von dieser Veränderlichkeit unabhängig sind. Man darf daher jede endliche Gleichung zwischen den Koordinaten \(x,y,z,r\), den Elementen und den anderen Größen, die unter der elliptischen Voraussetzung konstant sind und ebenfalls unter der Bezeichnung Elemente zusammengefasst werden können, differenzieren und dabei nur die Differenziale der Koordinaten berücksichtigen. Zieht man dieses Ergebnis vom vollständigen Differenzial ab, so bleibt das Differenzial der gegebenen Gleichung übrig, bei dem allein die Elemente als veränderlich behandelt werden. Offenbar darf also auch auf diese Weise differenziert werden. Dieser bemerkenswerte Satz kann bisweilen dazu dienen, die Rechnungen abzukürzen. Man sieht jedoch, dass man stets auch ohne ihn auskommen kann.
